\documentclass[11pt]{article}

\usepackage[margin=1in]{geometry}
\usepackage[T1]{fontenc}
\usepackage[utf8]{inputenc}
\usepackage{lmodern}
\usepackage{amsmath}
\usepackage{amssymb}
\usepackage{hyperref}
\usepackage{enumitem}
\usepackage{booktabs}
\usepackage{longtable}

\title{Internal Note: How Electrophysiology Ionic Models Are Coded in svMultiPhysics}
\author{}
\date{\today}

\begin{document}
\maketitle

\section{Purpose and Scope}

This note documents how the electrophysiology ionic models are implemented in the current C++ codebase. The focus is on the code architecture and data flow, not on the physiology of each model.

The main files involved are:
\begin{itemize}[leftmargin=*]
\item \texttt{Code/Source/solver/ionic\_model.h} and \texttt{ionic\_model.cpp}
\item \texttt{Code/Source/solver/cep\_ion.h} and \texttt{cep\_ion.cpp}
\item \texttt{Code/Source/solver/CepMod.h} and \texttt{CepMod.cpp}
\item \texttt{Code/Source/solver/read\_files.cpp}
\item \texttt{Code/Source/solver/Parameters.h} and \texttt{Parameters.cpp}
\item the concrete model files:
  \texttt{ionic\_aliev\_panfilov.*},
  \texttt{ionic\_bueno\_orovio.*},
  \texttt{ionic\_fitzhugh\_nagumo.*},
  \texttt{ionic\_ttp.*}
\end{itemize}

At the time of writing, the registered ionic models are:
\begin{itemize}[leftmargin=*]
\item Aliev--Panfilov (\texttt{AP})
\item Bueno--Orovio (\texttt{BO})
\item FitzHugh--Nagumo (\texttt{FN})
\item ten Tusscher--Panfilov 2006 (\texttt{TTP})
\end{itemize}

\section{High-Level Design}

The code separates electrophysiology into two layers:
\begin{enumerate}[leftmargin=*]
\item a \textbf{global CEP layer}, which stores nodal ionic unknowns, copies the solved transmembrane potential into the local ionic state, loops over all nodes, and manages domain averaging and optional electromechanics coupling;
\item a \textbf{local ionic-model layer}, which advances the ODE system for one node and one time step.
\end{enumerate}

The global layer lives mainly in \texttt{cep\_ion.cpp}. The local layer is represented by the abstract base class \texttt{IonicModel}.

Conceptually, the local model solves
\[
\frac{\partial v}{\partial t} + I_{\mathrm{ion}}(v,\mathbf{w}) = I_{\mathrm{ext}},
\qquad
\frac{\partial \mathbf{w}}{\partial t} = \mathbf{f}(v,\mathbf{w}),
\]
where \texttt{X} stores the non-gating state variables and \texttt{Xg} stores the gating variables.

\section{Core Data Structures}

\subsection{\texttt{IonicModel}}

\texttt{IonicModel} is the shared interface for all ionic models. It owns:
\begin{itemize}[leftmargin=*]
\item the initial values for state variables, \texttt{initial\_X};
\item the initial values for gating variables, \texttt{initial\_Xg};
\item voltage/time scaling constants:
  \texttt{Vrest}, \texttt{Vscale}, \texttt{Tscale}, and \texttt{Voffset}.
\end{itemize}

The important public methods are:
\begin{itemize}[leftmargin=*]
\item \texttt{init(X, Xg)}: fills local vectors with the model's initial conditions;
\item \texttt{integ(...)}: advances the model by one local ODE time step;
\item \texttt{read\_parameters(...)} and \texttt{distribute\_parameters(...)}: populate model parameters from parsed XML and broadcast them in parallel;
\item \texttt{get\_calcium\_index()}: returns the state index used as the calcium or calcium-proxy output for electromechanics coupling;
\item \texttt{get\_output\_variables()}: controls which ionic state variables can be exported.
\end{itemize}

The important protected virtual methods are:
\begin{itemize}[leftmargin=*]
\item \texttt{update\_g(...)}: update the gating variables analytically;
\item \texttt{getf(...)}: evaluate the right-hand side of the state equations;
\item \texttt{getj(...)}: evaluate the Jacobian used by the implicit Crank--Nicolson solver.
\end{itemize}

\subsection{\texttt{cepModelType} and \texttt{CepMod}}

\texttt{cepModelType} stores the domain-level setup for one electrophysiology region:
\begin{itemize}[leftmargin=*]
\item selected model type (\texttt{cepType});
\item counts \texttt{nX} and \texttt{nG};
\item myocardial zone id \texttt{imyo};
\item local ODE step size \texttt{dt};
\item stretch-activated-current coefficient \texttt{Ksac};
\item stimulus definition \texttt{Istim};
\item ODE integration options \texttt{odes};
\item the actual model instance, \texttt{ionic\_model}.
\end{itemize}

\texttt{CepMod} stores equation-wide CEP data. The most important field here is
\texttt{Xion}, an array of size
\[
\texttt{nXion} \times \texttt{tnNo},
\]
where \texttt{tnNo} is the total node count and \texttt{nXion} is the maximum number of ionic unknowns over all CEP domains. Each node column stores:
\[
[\text{state variables } X \; | \; \text{gating variables } X_g].
\]

This is a key implementation choice: every node uses the same global storage width, even if different domains use ionic models with different numbers of variables.

\section{How Models Are Selected and Constructed}

\subsection{String-to-model mapping}

The XML-facing model names are translated into the internal enum
\texttt{ElectrophysiologyModelType} in \texttt{CepMod.cpp}. For example:
\begin{itemize}[leftmargin=*]
\item \texttt{ap} or \texttt{aliev-panfilov} $\rightarrow$ \texttt{AP}
\item \texttt{bo} or \texttt{bueno-orovio} $\rightarrow$ \texttt{BO}
\item \texttt{fn} or \texttt{fitzhugh-nagumo} $\rightarrow$ \texttt{FN}
\item \texttt{ttp} or \texttt{tentusscher-panfilov} $\rightarrow$ \texttt{TTP}
\end{itemize}

\subsection{Factory registration}

Concrete ionic models self-register with \texttt{IonicModelFactory} using the macro
\begin{verbatim}
REGISTER_IONIC_MODEL("AP", AlievPanfilov);
\end{verbatim}
or the equivalent label for the other models.

This means the parser does not hard-code constructors for each ionic model. Instead, \texttt{read\_cep\_domain(...)} calls
\begin{verbatim}
IonicModelFactory::create_model(model_name)
\end{verbatim}
to build the domain's model object.

\section{How Parameters Are Parsed}

\subsection{Per-model parameter objects}

Each ionic model defines a nested \texttt{Parameters} class derived from \texttt{IonicModelParameters}. That object declares:
\begin{itemize}[leftmargin=*]
\item the XML element name for that model;
\item scalar parameters through \texttt{add\_parameter(label, default, required)};
\item vector parameters for zone-dependent data when needed;
\item the labels and defaults of the initial states and initial gating variables.
\end{itemize}

For example:
\begin{itemize}[leftmargin=*]
\item \texttt{AP} and \texttt{FN} use scalar parameters;
\item \texttt{BO} stores many parameters as length-3 vectors for epi/endo/myo;
\item \texttt{TTP} exposes a large number of scalar physical parameters.
\end{itemize}

\subsection{How the domain parser sees all models}

When a \texttt{DomainParameters} object is constructed, it iterates over the registered ionic models and asks each model for a parameter object through
\begin{verbatim}
model.get_parameters()
\end{verbatim}
The resulting objects are stored in
\begin{verbatim}
std::map<std::string, std::unique_ptr<IonicModelParameters>> ionic_models;
\end{verbatim}

As a result, the XML parser is ready to recognize sub-sections for every registered model, even before it knows which one the domain will actually use.

\subsection{What \texttt{read\_cep\_domain(...)} does}

The CEP domain reader in \texttt{read\_files.cpp} performs the following steps:
\begin{enumerate}[leftmargin=*]
\item read \texttt{<Electrophysiology\_model>};
\item instantiate the matching \texttt{IonicModel};
\item call \texttt{ionic\_model->read\_parameters(...)} using the model-specific parsed parameter object;
\item set \texttt{nX} and \texttt{nG} from the created model instance;
\item update the global \texttt{cep\_mod.nXion} if this domain has the largest ionic state count seen so far;
\item read conductivity, anisotropy, stimulus, ODE solver choice, local ODE time step, myocardial zone, and stretch feedback coefficient.
\end{enumerate}

\section{Global Runtime Flow}

\subsection{Initialization}

\texttt{cep\_ion::cep\_init(...)} initializes \texttt{CepMod::Xion}. For each CEP node:
\begin{enumerate}[leftmargin=*]
\item the relevant domain ionic model is asked to fill local vectors \texttt{Xl} and \texttt{Xgl};
\item these values are written into the node's column of \texttt{Xion};
\item if multiple CEP domains overlap a node, their ionic states are averaged.
\end{enumerate}

This averaging behavior is already present in the implementation and is important if multiple domain definitions contribute to one node.

\subsection{Per-time-step update}

\texttt{cep\_ion::cep\_integ(...)} performs one electrophysiology update after the diffusion solve. The main sequence is:
\begin{enumerate}[leftmargin=*]
\item optionally compute fiber stretch if electromechanics coupling is enabled;
\item copy the solved action potential from the global solution vector into \texttt{Xion(0,Ac)} for each node, except on the very first pass where \texttt{Xion} is already initialized;
\item for each CEP node and each applicable CEP domain, slice local vectors \texttt{Xl} and \texttt{Xgl} out of \texttt{Xion};
\item call \texttt{cep\_integ\_l(...)} to advance the local ionic ODE system;
\item write the updated local state back into \texttt{Xion};
\item copy the updated first state variable back into the global solution vector.
\end{enumerate}

Again, if several CEP domains contribute at a node, the updated states are averaged.

\subsection{Local ODE stepping}

\texttt{cep\_ion::cep\_integ\_l(...)} is the bridge between the global CEP equation and the local ionic model. It:
\begin{itemize}[leftmargin=*]
\item computes the stretch feedback coefficient
\[
K_{\mathrm{sac}} =
\begin{cases}
\texttt{cep.Ksac}\,(\sqrt{I4f}-1), & I4f > 1, \\
0, & \text{otherwise},
\end{cases}
\]
\item computes the number of local substeps as
\[
\texttt{nt} = \lfloor \texttt{dt} / \texttt{cep.dt} \rfloor,
\]
where the outer solver step is \texttt{dt} and the ionic-model step is \texttt{cep.dt};
\item constructs the stimulus window from \texttt{Ts}, \texttt{Td}, and \texttt{CL};
\item calls
\begin{verbatim}
cep.ionic_model->integ(...)
\end{verbatim}
once per local substep.
\end{itemize}

One practical implication is that ionic models support subcycling relative to the main solver time step.

\section{What Happens Inside \texttt{IonicModel::integ}}

The \texttt{integ(...)} method dispatches to one of three schemes:
\begin{itemize}[leftmargin=*]
\item Forward Euler (\texttt{FE})
\item classical fourth-order Runge--Kutta (\texttt{RK4})
\item Crank--Nicolson with Newton iteration (\texttt{CN2})
\end{itemize}

\subsection{State versus gating variables}

The code explicitly separates:
\begin{itemize}[leftmargin=*]
\item \texttt{X}: the main state variables, including the transmembrane potential at index 0;
\item \texttt{Xg}: the gating variables.
\end{itemize}

The integration logic assumes that \texttt{update\_g(...)} can advance \texttt{Xg} analytically or semi-analytically from the current state. The TTP model uses this pathway extensively. The AP, BO, and FN models have no separate gating-variable vector, so their \texttt{update\_g(...)} implementation is empty and all unknowns live in \texttt{X}.

\subsection{Scaling conventions}

Before evaluating the model equations, the base class may rescale voltage and time using the per-model constants:
\begin{itemize}[leftmargin=*]
\item \texttt{Vscale}
\item \texttt{Tscale}
\item \texttt{Voffset}
\end{itemize}

This is why some models can be coded in dimensionless form while still receiving dimensional inputs from the surrounding solver. For example:
\begin{itemize}[leftmargin=*]
\item AP uses \texttt{Vscale = 100.0}, \texttt{Tscale = 12.90}, \texttt{Voffset = -80.0};
\item BO uses \texttt{Vscale = 85.70}, \texttt{Tscale = 1.0}, \texttt{Voffset = -84.0};
\item FN and TTP are effectively used without such rescaling.
\end{itemize}

\subsection{Stretch-activated current}

The base class converts the scalar stretch coefficient into an additive current
\[
I_{\mathrm{sac}} = K_{\mathrm{sac}} \, (V_{\mathrm{rest}} - V).
\]
This current is scaled consistently with the model's voltage/time scaling and then passed to \texttt{getf(...)}.

\subsection{Implicit integration}

For \texttt{CN2}, the base class:
\begin{enumerate}[leftmargin=*]
\item evaluates the old-time residual;
\item performs Newton iterations on the state vector \texttt{X};
\item calls \texttt{getj(...)} for the Jacobian;
\item updates the gating variables only after the nonlinear solve.
\end{enumerate}

Models that do not override \texttt{getj(...)} will throw a \texttt{NotImplementedException}. In the current code, AP, BO, FN, and TTP all implement a Jacobian, but \texttt{read\_cep\_domain(...)} explicitly blocks implicit integration for TTP because it can give unexpected results.

\section{Patterns Used by the Concrete Models}

\subsection{Aliev--Panfilov and FitzHugh--Nagumo}

These two are the simplest implementations:
\begin{itemize}[leftmargin=*]
\item all unknowns are stored in \texttt{X};
\item \texttt{Xg} is empty;
\item \texttt{update\_g(...)} does nothing;
\item \texttt{getf(...)} and \texttt{getj(...)} are short closed-form expressions.
\end{itemize}

This is the easiest pattern to copy when adding a simple phenomenological model.

\subsection{Bueno--Orovio}

Bueno--Orovio is still coded with all unknowns in \texttt{X}, but its parameters are zone-dependent. The \texttt{zone\_id} passed into \texttt{getf(...)} and \texttt{getj(...)} is used as a 1-based selector for the epi/endo/myo parameter vectors.

This is an important coding convention:
\begin{itemize}[leftmargin=*]
\item \texttt{zone\_id == 1} means epicardium;
\item \texttt{zone\_id == 2} means endocardium or Purkinje;
\item \texttt{zone\_id == 3} means mid-myocardium.
\end{itemize}

\subsection{ten Tusscher--Panfilov}

TTP is the most detailed implementation and best illustrates the intended split between \texttt{X} and \texttt{Xg}:
\begin{itemize}[leftmargin=*]
\item \texttt{X} stores voltage, ionic concentrations, and the ryanodine-receptor state;
\item \texttt{Xg} stores gating variables for the channel kinetics;
\item \texttt{update\_g(...)} performs exact exponential gate updates;
\item \texttt{getf(...)} computes the current balance and calcium handling terms;
\item \texttt{getj(...)} provides the full Jacobian for the state equations.
\end{itemize}

This is the main example to follow for a biophysically richer model.

\section{Initial Conditions and Outputs}

\subsection{Initial conditions}

Every model declares its initial state labels and default values as static lists of
\begin{verbatim}
std::pair<std::string, double>
\end{verbatim}
These labels are used both for documentation and for XML parsing of model-specific
\texttt{<Initial\_states>} and \texttt{<Initial\_gating\_variables>} blocks.

\subsection{Exported ionic outputs}

By default, \texttt{IonicModel::get\_output\_variables()} exports one variable named \texttt{Calcium}, mapped to \texttt{get\_calcium\_index()}.

This is slightly overloaded in the phenomenological models:
\begin{itemize}[leftmargin=*]
\item in AP and FN, the ``calcium'' output is actually the recovery variable \texttt{w};
\item in BO, it is the slow inward gate \texttt{s};
\item in TTP, it is the actual intracellular calcium concentration \texttt{Ca\_i}.
\end{itemize}

The dynamic output registration happens in \texttt{read\_files.cpp} when the user requests spatial output of ionic state variables.

\section{How to Add a New Ionic Model}

To add a new model consistent with the current design:
\begin{enumerate}[leftmargin=*]
\item create \texttt{ionic\_<name>.h/.cpp} with a class derived from \texttt{IonicModel};
\item define the model label, initial state lists, and calcium index;
\item add a nested \texttt{Parameters} class derived from \texttt{IonicModelParameters};
\item override \texttt{get\_parameters()}, \texttt{read\_parameters()}, and \texttt{distribute\_parameters()};
\item implement \texttt{update\_g(...)} if the model uses separately stored gating variables;
\item implement \texttt{getf(...)} and, if implicit stepping is desired, \texttt{getj(...)};
\item register the model with \texttt{REGISTER\_IONIC\_MODEL};
\item extend the enum/string maps in \texttt{CepMod.h} and \texttt{CepMod.cpp};
\item ensure the model label is exposed to the XML parser through the factory-driven parameter creation path.
\end{enumerate}

In practice, the existing implementations suggest two common templates:
\begin{itemize}[leftmargin=*]
\item use AP or FN as the pattern for a simple phenomenological model;
\item use TTP as the pattern for a model with separate gating dynamics.
\end{itemize}

\section{Implementation Notes and Caveats}

\begin{itemize}[leftmargin=*]
\item The first entry of \texttt{X} is assumed to be the transmembrane potential throughout the CEP pipeline.
\item The code can average ionic states over overlapping CEP domains at a node.
\item The subcycling count is computed with integer truncation of \texttt{dt / cep.dt}. If the ratio is not an integer, the remainder is discarded.
\item The implicit solver checks absolute and relative residual norms, but reaching the maximum Newton iteration count does not itself trigger an error.
\item TTP has a Jacobian implementation, but the parser deliberately rejects implicit use of TTP.
\item The calcium output is a physiological calcium concentration only for some models; in reduced models it is a coupling proxy.
\end{itemize}

\section{Summary}

The ionic-model implementation is built around a small, fairly clean abstraction:
\begin{itemize}[leftmargin=*]
\item \texttt{cep\_ion.cpp} manages nodal storage and calls the local model;
\item \texttt{IonicModel} owns the common integration and scaling machinery;
\item concrete model classes provide only model-specific equations, parameters, and optional gate updates;
\item the factory and parameter system make new models pluggable with limited changes outside the model files and CEP enum maps.
\end{itemize}

That split is the main thing to keep in mind when navigating or extending the electrophysiology implementation in this codebase.

\end{document}
